% 復習5　問いを磨く（記入例）
\session{5}{1}{60分}{問いを磨く}{AIによる批判的検証}{}

\goals{
  \item AIを「厳しい査読者」として使い、問いの弱点を洗い出せる
  \item 指摘に対して、補強するか問いを修正するかを、根拠をもって決められる
  \item 問い・背景・先行研究・調べ方の見通しを1枚にまとめられる
}

\work{準備}{AIに提示する材料}
\abox[仮の問い]{14mm}{愛知県の中小企業では、なぜ生成AIの業務利用が進まないのか。}
\abox[背景（200字程度）]{28mm}{生成AIの業務利用は大企業を中心に広がっているが、中小企業での利用は限られていることが白書などで示されている。試験導入した中小製造業の研究では、費用よりも「使う業務を決められない」「社内で教える人がいない」ことが本格導入の障壁になっていた（復習4の文献）。しかし、地域の中小企業が実際にどの業務で使い始め、何に困っているかは、まだ十分にわかっていない。}
\abox[主な先行研究（3件）]{20mm}{① 復習4で精読した文献（試験導入した中小製造業の聞き取り調査）　② 中小企業庁「中小企業白書」　③ 総務省「情報通信白書」}

\work[70mm]{ワーク1・2}{査読者の指摘に、補強か修正で答える}
\begin{promptbox}[査読者として批判させる]
あなたは厳しい査読者です。次のリサーチ・クエスチョンについて、1. 数か月で調べて答えられるか　2. 先行研究に対する新しさがあるか　3. 範囲は適切か　4. 用語は明確か　の観点から問題点を指摘し、それぞれ改善案を示してください。問い：（あなたの問い）　背景：（200字程度）　主な先行研究：（3件）
\end{promptbox}
\begin{wstabh}{colspec={Q[l,wd=22mm,m] X[3,l,m] Q[c,wd=14mm,m] X[3,l,m]},row{2-Z}={ht=17mm}}
  観点 & AIの指摘（要点） & 補強／\newline 修正 & 自分の対応と、その根拠 \\
  答えられるか & \af{「なぜ進まないか」には、導入しない多くの企業を調べる必要があり、数か月では難しい} & \ans{修正} & \af{試験導入した企業に絞り、「どの業務で使い、何を障壁と感じているか」を聞く問いにする} \\
  新しさ & \af{先行研究（12社の聞き取り）と同じことを調べるだけにならないか} & \ans{補強} & \af{先行研究は「本格導入に進むか」を分ける要因。自分は「使い始めた業務」に注目する点が新しい} \\
  範囲 & \af{「中小企業」は業種が多すぎる} & \ans{修正} & \af{先行研究と比べられるよう、製造業に絞る} \\
  用語の明確さ & \af{「業務利用」が曖昧} & \ans{修正} & \af{「会社として、特定の業務に生成AIを使っていること」と定義し、社員の個人的な利用は含めない} \\
\end{wstabh}

\begin{promptbox}[反論に答える練習]
私の問いに対して「（AIが挙げた反論）」という批判があります。この批判に答える論点を3つ挙げ、それぞれどんな根拠が必要かを示してください。根拠が見つかっていない論点は、そのように明示してください。
\end{promptbox}

\work{修正の記録}{問いはどう変わったか}
\atwoboxes{修正前の問い}{修正後の問い（確定）}{24mm}
{愛知県の中小企業では、なぜ生成AIの業務利用が進まないのか。}
{生成AIを試験導入した愛知県の中小製造業は、どの業務で使い、本格導入の障壁を何と捉えているのか。}

\work[90mm]{ワーク3}{確定した問いを1枚にまとめる}
\begin{wstab}{colspec={Q[l,wd=30mm,m] X[l,m]},row{1}={ht=14mm},row{2-Z}={ht=22mm}}
  \hc{問い} & \af{生成AIを試験導入した愛知県の中小製造業は、どの業務で使い、本格導入の障壁を何と捉えているのか。} \\
  \hc{背景と意義} & \af{生成AIの業務利用は中小企業で限られている。試験導入から先に進めない理由がわかれば、これから導入する企業や、支援する自治体・金融機関にとって手がかりになる。} \\
  \hc{先行研究\newline でわかって\newline いること} & \af{試験導入した中小製造業では、費用より「使う業務を決められない」「教える人がいない」が障壁になる（復習4の文献）。白書から、企業規模による利用の差がわかる。使い始めた業務の具体像は未解決。} \\
  \hc{調べ方の見通し} & \af{① 白書・統計で全体の傾向を確認する　② 試験導入した中小製造業（アルバイト先と、その紹介で3〜5社）に聞き取りをする　③ 使っている業務と障壁を表に整理し、先行研究と比べる} \\
\end{wstab}

\work[40mm]{振り返り}{復習1〜5のチェック}
\begin{achecks}
  \item AIが挙げた文献・数値を、そのまま使わずに確かめた
  \item 問いの候補から、自分で基準を決めて選んだ
  \item 文献リストの必須項目をすべて揃えた
  \item AIの要約を原文と照らし合わせた
  \item AIの指摘を受けて、問いをどう変えたかを記録した
  \item AI活用記録を毎回書いた
\end{achecks}
\aimemoA{問いへの批判（査読者役）と、反論に答える論点の整理}
{「あなたは厳しい査読者です。4つの観点から問題点と改善案を」「この批判に答える論点を3つ、必要な根拠とともに」}
{4つの指摘のうち3つで問いを修正し、新しさの指摘には先行研究との違いを示して反論した。改善案の文言はそのまま使わず、自分で書き直した。}
{修正後の問いが、手元の文献と聞き取りで答えられるかを、先行研究の調査方法と照らして確かめた}

\begin{point}
  \item 修正前と修正後の問いを並べて記録することが最も大切。AIの指摘を受けて何をどう変えたかが、AI活用記録の中心になる。
  \item すべての指摘を受け入れる必要はない。根拠を示して反論できた点（記入例では「新しさ」）は、自分の判断の跡として残す。
  \item 「なぜ〜しないのか」のように、しない理由を問う問いは、調べる対象が大きくなりがち。記入例のように「した人（企業）」に絞ると答えやすくなる。
  \item 1枚のまとめは、そのままレポートの序論の骨組みになる。
\end{point}
